\begin{enumerate}

\item Identificar si la sección está abierta, cerrada unicelular o bicelular; dibujar su geometría y espesores.

\item Para la sección abierta emplear J \ensuremath{\approx} \ensuremath{\Sigma} b\_i t\_i\textsuperscript{3}/3 y tau\_max \ensuremath{\approx} M\_x t/J; NO se usa el flujo circulante de Bredt.

\item Comprobar que el par exterior coincide con el que transmiten los flujos, si la sección es cerrada.

\item Calcular el giro total con theta = M\_x L/(G J), o con la compatibilidad de las celdas.

\item Las magnitudes de la solución se han contrastado mediante cálculo analítico y el motor estructural de Studium cuando éste ofrece esa comprobación.

\end{enumerate}

\noindent\textbf{Operaciones y resultados intermedios.}

\begin{enumerate}

\item Constante abierta de Saint-Venant (suma de tiras).

\[\fitmath{y_{1}=\frac{\left(100+70\right)\cdot {2}^{3}}{3}\approx 453.3333\,\mathrm{mm^{4}}}\]

\item Giro total entre extremos.

\[\fitmath{y_{2}=\frac{4000\cdot 300}{26955\cdot \frac{\left(100+70\right)\cdot {2}^{3}}{3}}\approx 0.09820289\,\mathrm{rad}}\]

\item Tensión máxima aproximada en tira delgada.

\[\fitmath{y_{3}=\frac{4000\cdot 2}{\frac{\left(100+70\right)\cdot {2}^{3}}{3}}\approx 17.64706\,\mathrm{MPa}}\]

\end{enumerate}
